\chapter{Prompts utilizados para inferencia con LLMs}
\label{anexo:prompts_llm}

\section{Prompt de clasificación binaria: ansiedad}

\begin{lstlisting}[style=promptstyle]
Vas a recibir un texto que corresponde a una serie de mensajes escritos por una 
persona en una red social.
El texto puede o no reflejar indicios de ansiedad.

La ansiedad se asocia a una preocupación persistente y difícil de controlar, 
acompañada de nerviosismo, inquietud interna y una sensación continua de amenaza 
o anticipación negativa. Estas experiencias suelen generar malestar emocional 
sostenido y dificultad para relajarse.

Tu tarea es clasificar el texto únicamente en una de las siguientes categorías:

1: el texto muestra indicios de ansiedad
0: el texto no muestra indicios de ansiedad

El contenido de los mensajes se proporciona a continuación entre comillas triples:

"""{text}"""

Responde solo con el número correspondiente.
\end{lstlisting}


\section{Prompt de clasificación binaria: depresión}

\begin{lstlisting}[style=promptstyle]
Vas a recibir un texto que corresponde a una serie de mensajes escritos por una 
persona en una red social.
El texto puede o no reflejar indicios de depresión.

La depresión se asocia a un estado persistente de tristeza, apatía o vacío 
emocional, acompañado de una disminución del interés o la capacidad de disfrute en 
actividades habituales. Estas experiencias pueden incluir pensamientos negativos 
recurrentes, fatiga emocional y una percepción pesimista de uno mismo o del futuro.

Tu tarea es clasificar el texto únicamente en una de las siguientes categorías:

1: el texto muestra indicios de depresión
0: el texto no muestra indicios de depresión

El contenido de los mensajes se proporciona a continuación entre comillas triples:

"""{text}"""

Responde solo con el número correspondiente.
\end{lstlisting}


\section{Prompt de clasificación binaria: TCA}

\begin{lstlisting}[style=promptstyle]
Vas a recibir un texto que corresponde a una serie de mensajes escritos por una 
persona en una red social.
El texto puede o no reflejar indicios de un trastorno de la conducta alimentaria.

Los trastornos de la conducta alimentaria se asocian a una preocupación persistente 
por el peso corporal, la alimentación o la imagen corporal, que puede manifestarse 
mediante pensamientos rígidos o negativos sobre la comida, el cuerpo o el control 
del peso. Estas preocupaciones suelen generar malestar emocional, culpa, ansiedad 
en torno a la alimentación y conductas de evitación o control.

Tu tarea es clasificar el texto únicamente en una de las siguientes categorías:

1: el texto muestra indicios de un trastorno de la conducta alimentaria
0: el texto no muestra indicios de un trastorno de la conducta alimentaria

El contenido de los mensajes se proporciona a continuación entre comillas triples:

"""{text}"""

Responde solo con el número correspondiente.
\end{lstlisting}


\section{Prompt de clasificación multiclase}

\begin{lstlisting}[style=promptstyle]
Vas a recibir un texto que corresponde a una serie de mensajes escritos por una 
persona en una red social.
El texto puede reflejar indicios de uno de los siguientes estados psicológicos, 
o no reflejar ninguno de ellos.

Ansiedad: se asocia a una preocupación persistente y difícil de controlar, 
acompañada de nerviosismo, inquietud interna y una sensación continua de amenaza 
o anticipación negativa, generando malestar emocional y dificultad para relajarse.

Depresión: se asocia a un estado persistente de tristeza, apatía o vacío emocional, 
junto con una disminución del interés o la capacidad de disfrute en actividades 
habituales, pensamientos negativos recurrentes y una visión pesimista de uno mismo 
o del futuro.

Trastornos de la conducta alimentaria: se asocian a una preocupación persistente 
por el peso, la alimentación o la imagen corporal, que puede manifestarse mediante 
pensamientos rígidos o negativos sobre la comida o el cuerpo, así como malestar 
emocional, culpa o ansiedad en torno a la alimentación.

Tu tarea es clasificar el texto únicamente en una de las siguientes categorías:

0: el texto no muestra indicios de ninguno de los estados descritos
1: el texto muestra indicios de ansiedad
2: el texto muestra indicios de depresión
3: el texto muestra indicios de un trastorno de la conducta alimentaria

El contenido de los mensajes se proporciona a continuación entre comillas triples:

"""{text}"""

Responde solo con el número correspondiente.
\end{lstlisting}
